\chapter{The uncertainty head and the Laplace likelihood}\label{ch:uncertainty}

\begin{onething}
A small extra head can learn, for every pixel, how large an error the model expects there, by minimising a
Laplace negative log-likelihood whose best answer is exactly that expected error; but when the residual is
divided by this prediction, real targets fade, because a head that predicts its own error well also predicts
large errors on anomalies whenever it can see any trace of them.
\end{onething}

This chapter is history told as a lesson. It covers one idea, the \emph{uncertainty head}, from its proposal
on 2026-10-19 to its failure on 2026-10-20 as the fine-tuning run FT06. Along the way you meet three pieces of
mathematics that are worth owning for their own sake: the Laplace distribution, the method of maximum
likelihood applied to a scale parameter, and Spearman's rank correlation. The chapter then shows, with
arithmetic you can do on paper, why a head that passed its own sanity check still made detection worse.

A note on the provenance tags in this chapter and the next two. To keep lines short they abbreviate
paths:
\begin{itemize}
\item run folders such as \file{FT06-uncertainty/}: in\\ \file{research/experiments/indradhanu-finetune/};
\item \file{bench-resolution/}, \file{scoring-gains/} and \file{decoder-capacity/}: in\\
  \file{research/experiments/};
\item \file{finetune_trainer.py}: in \file{app/foundation_models/trainers/};
\item \file{scorers.py}: in \file{app/evaluation/};
\item \file{uncertainty-head.md} is in \file{docs/proposals/}, and \file{lld/finetune.md} in \file{docs/}.
\end{itemize}

\section{Where we were, and the gap the score leaves}\label{sec:uncertainty:where}

\Cref{ch:shipping} shipped FT04 as Indradhanu V2. Its detection tied the older v0.2.0 model, and on the same
day an experiment tried to win detection back on the \emph{score} side while holding FT04's reconstruction
fixed: blurring the residual, fusing residual Mahalanobis with PCA-RX, and whitening each background type
separately (\cref{ch:scoringgains}). None of them beat plain residual Mahalanobis on real objects caught. The
experiment's verdict was that ``the gains question goes to the model side''
\tagM{scoring-gains/README.md}.

Recall from \cref{ch:residualmaha} what residual Mahalanobis (RM) does. It takes the residual
$\vect{r}(p)=\vect{x}(p)-\xhat(p)$ at each pixel $p$, fits \emph{one} mean and \emph{one} covariance matrix for
the whole scene, and scores each pixel by its Mahalanobis distance under that covariance. The covariance
learns which \emph{directions} of error the scene shares: a haze that lifts the blue bands everywhere, a
detector stripe, dark water. What it cannot learn is which \emph{places} are hard. The model's errors are not
alike from place to place. Field edges, fine texture, sun glint, thin cloud and shorelines are hard to
reconstruct everywhere, and their large residuals are exactly RM's false alarms. RM has no way to know that
``the model always struggles here''.

\begin{analogy}
A weather forecaster who says ``20\,mm of rain, give or take 2'' and one who says ``20\,mm, give or take 15''
can both see 35\,mm fall. For the first it is a shock; for the second it is a wet day. What makes an
observation surprising is the miss \emph{divided by} how big a miss the forecaster expected. The uncertainty
head is a forecaster of the reconstruction's own error. Later in the chapter the analogy acquires a twist:
a forecaster who is never told about your town, only about the towns around it, will still see a storm
parked over your town spilling rain into every neighbour, and will say ``something big near there, give or
take a lot''. The storm then stops looking surprising.
\end{analogy}

The obvious fix is to give each pixel its own ``give or take'': a number $b(p)>0$, the error size the model
expects at $p$, and to score $\vect{r}(p)/b(p)$ instead of $\vect{r}(p)$. Two pixels with the same residual
$0.3$ then differ: at a field edge where the head expects $b=0.3$ the ratio is $0.3/0.3=1.0$, ordinary; on a
flat field where it expects $b=0.05$ the ratio is $0.3/0.05=6.0$, a surprise. The rest of this chapter is about
how to learn $b(p)$ honestly, and why it still did not work.

\section{A probability model for errors: the Laplace distribution}\label{sec:uncertainty:laplace}

To learn $b$ by optimisation we need a loss whose minimum is ``the right size of error''. The cleanest way to
build one is to write down a probability model for the error and ask which $b$ makes the observed errors most
probable. The model chosen is the Laplace distribution, because it is the distribution that matches the L1
reconstruction loss of \cref{ch:losses}, as you will see in \cref{sec:uncertainty:nll}.

\begin{definition}[Probability density]
A \emph{probability density}\index{probability density} $f(r)$ for a real-valued quantity $r$ is a function
with $f(r)\ge 0$ everywhere and $\int_{-\infty}^{\infty} f(r)\,dr=1$. The probability that $r$ lands in a
small interval of width $\Delta$ around $r_0$ is approximately $f(r_0)\,\Delta$.
\end{definition}

\begin{definition}[Laplace distribution]
The \emph{Laplace distribution}\index{Laplace distribution} with centre $0$ and \emph{scale}\index{scale
parameter} $b>0$ has density
\begin{equation}\label{eq:uncertainty:laplace}
  f(r\mid b) \;=\; \frac{1}{2b}\,\exp\!\left(-\frac{\abs{r}}{b}\right).
\end{equation}
\end{definition}

Read \cref{eq:uncertainty:laplace} in words: the density is highest at $r=0$, where it equals $1/(2b)$, and
falls away exponentially on both sides as the size of the error $\abs{r}$ grows, at a rate set by $b$. A small
$b$ means a tall narrow peak (errors are expected to be small); a large $b$ a low wide one.

\begin{figure}[htbp]
\centering
\begin{tikzpicture}
\begin{axis}[width=0.8\linewidth,height=5.5cm,xlabel={error $r$},ylabel={density $f(r\mid b)$},
  domain=-3:3,samples=241,xmin=-3,xmax=3,ymin=0,ymax=1.05,legend pos=north east,legend cell align=left,
  grid=major,grid style={gray!20}]
\addplot[spec,thick] {exp(-abs(x)/0.5)/(2*0.5)};
\addlegendentry{$b=0.5$}
\addplot[spat,thick,dashed] {exp(-abs(x)/1)/(2*1)};
\addlegendentry{$b=1$}
\addplot[loss,thick,dotted] {exp(-abs(x)/2)/(2*2)};
\addlegendentry{$b=2$}
\addplot[maskgrey,thin] {exp(-x^2/2)/sqrt(2*pi)};
\addlegendentry{Gaussian, $\sigma=1$}
\end{axis}
\end{tikzpicture}
\caption{The Laplace density for three scales, drawn to scale from \cref{eq:uncertainty:laplace}, with a
standard Gaussian in grey for comparison. The Laplace has a sharp peak at zero and heavier tails: a large
error is less surprising under a Laplace than under a Gaussian of similar width.}
\label{fig:uncertainty:laplace}
\end{figure}

\subsection{It is a density: the integral is one}

You should check that \cref{eq:uncertainty:laplace} really integrates to one. The function is symmetric
($f(-r)=f(r)$, since $\abs{-r}=\abs{r}$), so the integral is twice the integral over $r\ge 0$, where
$\abs{r}=r$:
\begin{align*}
\int_{-\infty}^{\infty}\frac{1}{2b}e^{-\abs{r}/b}\,dr
  &= 2\int_{0}^{\infty}\frac{1}{2b}e^{-r/b}\,dr
  \intertext{cancel the 2 against the $\tfrac12$:}
  &= \frac{1}{b}\int_{0}^{\infty}e^{-r/b}\,dr
  \intertext{an antiderivative of $e^{-r/b}$ is $-b\,e^{-r/b}$ (differentiate it to check: $-b\cdot(-\tfrac1b)e^{-r/b}=e^{-r/b}$):}
  &= \frac{1}{b}\Big[-b\,e^{-r/b}\Big]_{0}^{\infty}
  \intertext{at $r\to\infty$ the exponential is $0$; at $r=0$ it is $1$:}
  &= \frac{1}{b}\big(0-(-b)\big) \;=\; \frac{1}{b}\cdot b \;=\; 1.
\end{align*}

\subsection{The scale is the mean absolute error}

The single most useful fact about the Laplace scale is that it equals the average size of the error,
$\E\abs{r}=b$. Derive it the same way, using the symmetry and then integration by parts (the rule
$\int u\,dv = uv-\int v\,du$):
\begin{align*}
\E\abs{r} &= \int_{-\infty}^{\infty}\abs{r}\,\frac{1}{2b}e^{-\abs{r}/b}\,dr
  \;=\; \frac{1}{b}\int_{0}^{\infty} r\,e^{-r/b}\,dr
  \intertext{take $u=r$ (so $du=dr$) and $dv=e^{-r/b}dr$ (so $v=-b\,e^{-r/b}$):}
  &= \frac{1}{b}\left(\Big[-b\,r\,e^{-r/b}\Big]_0^{\infty} + b\int_0^{\infty}e^{-r/b}\,dr\right)
  \intertext{the bracket is $0$ at both ends ($r=0$ kills it at the bottom; the exponential beats $r$ at the top); the integral was found above to be $b$:}
  &= \frac{1}{b}\left(0 + b\cdot b\right) \;=\; b.
\end{align*}
So when the head outputs $b(p)$ it is literally predicting ``the mean absolute error I expect at this
pixel''. That is why the proposal could describe $b$ in plain words as the expected error size.

\section{Likelihood, and the negative log-likelihood}\label{sec:uncertainty:nll}

\begin{definition}[Likelihood]
Given an observed value $r$ and a model with parameter $b$, the \emph{likelihood}\index{likelihood} of $b$ is
the density of the observation viewed as a function of the parameter: $\mathcal{L}(b)=f(r\mid b)$. The
\emph{maximum-likelihood}\index{maximum likelihood} estimate of $b$ is the value that makes $\mathcal{L}(b)$
largest.
\end{definition}

Maximising a product of small numbers is awkward, so we take logarithms (which turns products into sums and
does not move the maximum, because $\log$ is increasing) and flip the sign (so that maximising becomes
minimising, which is what an optimiser does). The result is the \emph{negative log-likelihood}\index{negative
log-likelihood} (NLL). For the Laplace:
\begin{align}
-\log f(r\mid b) &= -\log\!\left(\frac{1}{2b}\exp\!\left(-\frac{\abs{r}}{b}\right)\right)
  \nonumber\\
  \intertext{the log of a product is the sum of the logs:}
  &= -\log\frac{1}{2b} \;-\; \log\exp\!\left(-\frac{\abs{r}}{b}\right) \nonumber\\
  \intertext{$\log(1/y)=-\log y$, and $\log\exp(z)=z$:}
  &= \log(2b) + \frac{\abs{r}}{b} \nonumber\\
  \intertext{and $\log(2b)=\log 2+\log b$:}
  &= \frac{\abs{r}}{b} + \log b + \log 2. \label{eq:uncertainty:nll}
\end{align}
The last term, $\log 2\approx 0.693$, does not depend on $b$, so it cannot change which $b$ is best. The code
drops it. What is left is the per-pixel loss the trainer minimises:
\begin{equation}\label{eq:uncertainty:loss}
  \ell(b;\,\abs{r}) \;=\; \frac{\abs{r}}{b} + \log b .
\end{equation}

Two forces pull on $b$. The first term, $\abs{r}/b$, punishes a $b$ that is too small: claim to be confident
and then miss, and you pay a lot. The second term, $\log b$, punishes a $b$ that is too large: you cannot
escape the first term by claiming to know nothing, because vagueness costs $\log b$.

\begin{keypoint}
The Laplace NLL is the L1 error divided by a predicted scale, plus the log of that scale; with $b$ fixed it is
the L1 loss of \cref{ch:losses} up to a constant factor and offset, which is why the Laplace is the
distribution that ``matches'' L1.
\end{keypoint}

\begin{worked}[label={we:uncertainty:three}]{The loss at three scales}
Take one pixel whose error is $\abs{r}=0.2$ and try $b=0.1$, $0.2$ and $0.4$. Do it on paper with a
calculator for the natural logarithm.
\begin{align*}
b=0.1:&\quad \frac{0.2}{0.1}+\ln 0.1 = 2.0 + (-2.303) = -0.303\\
b=0.2:&\quad \frac{0.2}{0.2}+\ln 0.2 = 1.0 + (-1.609) = -0.609\\
b=0.4:&\quad \frac{0.2}{0.4}+\ln 0.4 = 0.5 + (-0.916) = -0.416
\end{align*}
The lowest of the three is at $b=0.2$, the error itself. At $b=0.1$ the first term is large (2.0); at
$b=0.4$ the first term is small (0.5) but the second has risen from $-1.609$ to $-0.916$. Negative losses are normal here: a density can exceed 1, so its log can be positive and the NLL negative.
\end{worked}

\begin{figure}[htbp]
\centering
\begin{tikzpicture}
\begin{axis}[width=0.75\linewidth,height=5.5cm,xlabel={predicted scale $b$},ylabel={$\ell(b)=0.2/b+\ln b$},
  domain=0.06:1.0,samples=200,xmin=0,xmax=1.0,ymin=-0.8,ymax=1.0,grid=major,grid style={gray!20}]
\addplot[loss,thick] {0.2/x+ln(x)};
\addplot[only marks,mark=*,attn] coordinates {(0.1,-0.303) (0.2,-0.609) (0.4,-0.416)};
\draw[dashed,spec] (axis cs:0.2,-0.8) -- (axis cs:0.2,-0.609);
\node[spec,anchor=north west,font=\small] at (axis cs:0.21,-0.62) {minimum at $b=\abs{r}=0.2$};
\end{axis}
\end{tikzpicture}
\caption{The per-pixel loss of \cref{eq:uncertainty:loss} for $\abs{r}=0.2$, drawn to scale, with the three
points of \cref{we:uncertainty:three}. The curve rises steeply to the left (overconfidence) and slowly to the
right (vagueness).}
\label{fig:uncertainty:losscurve}
\end{figure}

\section{The best scale is the error: minimising over \texorpdfstring{$b$}{b}}\label{sec:uncertainty:min}

The worked example suggests that the minimum sits at $b=\abs{r}$. Prove it with calculus. At a minimum of a
smooth function the derivative is zero.
\begin{align*}
\frac{d\ell}{db} &= \frac{d}{db}\left(\abs{r}\,b^{-1}\right) + \frac{d}{db}\log b
  \intertext{the power rule gives $\frac{d}{db}b^{-1}=-b^{-2}$; and $\frac{d}{db}\log b = \frac1b$:}
  &= -\frac{\abs{r}}{b^2} + \frac{1}{b}.
  \intertext{Set it to zero:}
  0 &= -\frac{\abs{r}}{b^2} + \frac{1}{b}
  \intertext{multiply both sides by $b^2$ (allowed, because $b>0$):}
  0 &= -\abs{r} + b
  \intertext{so}
  b^\star &= \abs{r}.
\end{align*}
Is it a minimum and not a maximum? Differentiate again:
\begin{align*}
\frac{d^2\ell}{db^2} &= \frac{2\abs{r}}{b^3} - \frac{1}{b^2},
  \intertext{and at $b=\abs{r}$:}
\frac{d^2\ell}{db^2}\Big|_{b=\abs{r}} &= \frac{2\abs{r}}{\abs{r}^3}-\frac{1}{\abs{r}^2}
  = \frac{2}{\abs{r}^2}-\frac{1}{\abs{r}^2} = \frac{1}{\abs{r}^2} > 0 .
\end{align*}
A positive second derivative means the curve bends upward there: a minimum.

\begin{worked}[label={we:uncertainty:deriv}]{The derivative at three points}
Keep $\abs{r}=0.2$. Evaluate $d\ell/db=-\abs{r}/b^2+1/b$ at $b=0.1$, $0.2$, $0.4$.
\begin{align*}
b=0.1:&\quad -\frac{0.2}{0.01}+\frac{1}{0.1} = -20 + 10 = -10 \quad\text{(negative: increasing $b$ lowers the loss)}\\
b=0.2:&\quad -\frac{0.2}{0.04}+\frac{1}{0.2} = -5 + 5 = 0 \quad\text{(flat: the minimum)}\\
b=0.4:&\quad -\frac{0.2}{0.16}+\frac{1}{0.4} = -1.25 + 2.5 = 1.25 \quad\text{(positive: decreasing $b$ lowers the loss)}
\end{align*}
Notice the asymmetry: at half the right answer the slope is $-10$, at twice it only $+1.25$. Gradient descent
pushes hard against overconfidence and gently against vagueness.
\end{worked}

\subsection{Many pixels sharing one scale: the mean absolute error}

A single pixel's error says nothing about ``how hard this kind of place is''. The head sees context, not the
pixel, so it cannot tell apart many pixels whose contexts look alike. Suppose $N$ such pixels, with errors
$\abs{r_1},\dots,\abs{r_N}$, must share one predicted $b$. Their average loss is
\begin{align*}
\bar\ell(b) &= \frac1N\sum_{i=1}^N\left(\frac{\abs{r_i}}{b}+\log b\right)
  \intertext{the $\log b$ is the same in every term, so it comes out of the average; the $1/b$ factors out of the sum:}
  &= \frac{1}{b}\cdot\frac1N\sum_{i=1}^{N}\abs{r_i} \;+\; \log b
   \;=\; \frac{m}{b}+\log b, \qquad m=\frac1N\sum_i\abs{r_i}.
\end{align*}
This is the single-pixel loss with $\abs{r}$ replaced by $m$, so by the same derivative $b^\star=m$: \textbf{the
best shared scale is the mean absolute error of the pixels that share it.} With infinitely many pixels, the
mean becomes the expectation, $b^\star(p)=\E\abs{r(p)}$, the typical error at that kind of place. This matches
the result $\E\abs{r}=b$ of \cref{sec:uncertainty:laplace}: if errors really were Laplace with scale $b$, the
maximum-likelihood fit recovers it.

\begin{worked}[label={we:uncertainty:four}]{Four pixels, one scale}
Four hidden pixels of a field edge have errors $0.10$, $0.30$, $0.20$, $0.40$. Find the best shared $b$ and
check that it beats a neighbour value. Do it on paper.
\begin{align*}
\text{sum of errors} &= 0.10+0.30+0.20+0.40 = 1.00\\
m &= 1.00/4 = 0.25\\
\bar\ell(0.25) &= 0.25/0.25 + \ln 0.25 = 1.000 + (-1.386) = -0.386\\
\bar\ell(0.20) &= 0.25/0.20 + \ln 0.20 = 1.250 + (-1.609) = -0.359\\
\bar\ell(0.30) &= 0.25/0.30 + \ln 0.30 = 0.833 + (-1.204) = -0.371
\end{align*}
$-0.386$ is the lowest: the mean absolute error wins against both neighbours.
\end{worked}

\section{Predicting \texorpdfstring{$\log b$}{log b}, and the detached residual}\label{sec:uncertainty:logb}

A network's last layer is a convolution, which can output any real number, including negative ones. A scale
must be positive. The standard trick, used here, is to let the network output $s=\log b$ and compute
$b=e^{s}$, which is positive for every real $s$. Rewrite the loss in terms of $s$:
\begin{align*}
\ell(s) &= \frac{\abs{r}}{e^{s}} + s = \abs{r}\,e^{-s} + s,
\intertext{and differentiate with respect to $s$ (the chain rule: $\frac{d}{ds}e^{-s}=e^{-s}\cdot(-1)$):}
\frac{d\ell}{ds} &= -\abs{r}\,e^{-s} + 1 .
\intertext{Setting it to zero gives $e^{-s}=1/\abs{r}$, so}
s^\star &= \log\abs{r}, \quad\text{that is } b^\star=e^{s^\star}=\abs{r}.
\end{align*}
The minimum is in the same place, as it must be: reparameterising does not move the best answer, only the
path to it.

\begin{worked}[label={we:uncertainty:chain}]{Gradients with respect to $b$ and to $\log b$}
This is the chain rule on a tiny example. Keep $\abs{r}=0.2$ and stand at $b=0.1$, so $s=\ln 0.1=-2.303$.
\begin{align*}
\frac{d\ell}{ds} &= -0.2\,e^{2.303}+1 = -0.2\times 10.0 + 1 = -2.0 + 1 = -1.0\\
\frac{d\ell}{db} &= -10 \quad\text{(from \cref{we:uncertainty:deriv})}\\
\frac{db}{ds} &= e^{s} = b = 0.1\\
\frac{d\ell}{ds} &= \frac{d\ell}{db}\cdot\frac{db}{ds} = -10\times 0.1 = -1.0 \quad\text{(the same, as the chain rule promises)}
\end{align*}
In $b$ the slope was $-10$; in $\log b$ it is $-1$. Near zero, where $b$ is small, the $\log b$ parameterisation
tames gradients that would otherwise explode as $1/b^2$.
\end{worked}

\paragraph{The clamp.} Before using the head's output, the trainer clamps it to the range
$-10\le s\le 5$ \tagM{finetune_trainer.py}. In terms of $b$ that is
$e^{-10}\approx 0.0000454$ to $e^{5}\approx 148.4$: wide enough never to bind on a sensible error, narrow
enough that one bad pixel cannot produce an infinite loss.

\paragraph{The detached residual.} The error $\abs{r}$ in the loss is computed inside \code{torch.no_grad()}:
it is \emph{detached}\index{detach}, meaning no gradient flows from the uncertainty loss back into the
reconstruction. Without this, the network could lower $\abs{r}/b+\log b$ by changing its \emph{errors} to suit
$b$, rather than changing $b$ to fit its errors. The docstring of \code{_scale_loss} says it in one line:
``detached: the target, not something to shrink''.

The trainer's loss, from that docstring:
\begin{lstlisting}
e(p) = mean over valid bands of |normalised residual|   # (B, H, W), computed under no_grad
loss = mean over hidden p of  e(p) / b(p) + log b(p)    # lowest where b(p) = e(p)
\end{lstlisting}
Note that $e(p)$ is a per-pixel mean over bands, so $b(p)$ is one number per pixel, not one per band. The
proposal rejected per-band scales on purpose: the per-band shape is RM's job.

\section{The head itself, counted by hand}\label{sec:uncertainty:head}

The \emph{uncertainty head}\index{uncertainty head} is a third decoder on the shared encoder features, beside
the reconstruction decoder and the segmentation head (\cref{ch:seghead}). It is a \code{SegFormerDecoder}
(\cref{ch:decoder}) built with width $64$ and one output channel:
\begin{lstlisting}
scale_head = SegFormerDecoder(embed_dims=[32, 64, 160, 256], decoder_dim=64, out_channels=1)
#   features F1..F4: (B,32,H/4,W/4) (B,64,H/8,W/8) (B,160,H/16,W/16) (B,256,H/32,W/32)
#   -> 1x1 projections to 64, bilinear up to H/4, concat     (B, 256, H/4, W/4)
#   -> 1x1 fuse + GELU                                       (B, 64,  H/4, W/4)
#   -> 3x3 refine + GELU                                     (B, 64,  H/4, W/4)
#   -> 3x3 conv to 1 x 16 sub-pixel channels                 (B, 16,  H/4, W/4)
#   -> PixelShuffle(4)                                       (B, 1,   H,   W)    = log b
\end{lstlisting}
Every convolution has $c_\text{in}\times c_\text{out}\times k\times k$ weights plus $c_\text{out}$ biases
(\cref{ch:shapes}).

\begin{worked}[label={we:uncertainty:count}]{The uncertainty head's 95,632 parameters}
Count each layer on paper, then add.
\begin{align*}
\text{projection from }32 &: 32\times 64 + 64 = 2{,}048+64 = 2{,}112\\
\text{projection from }64 &: 64\times 64 + 64 = 4{,}096+64 = 4{,}160\\
\text{projection from }160 &: 160\times 64 + 64 = 10{,}240+64 = 10{,}304\\
\text{projection from }256 &: 256\times 64 + 64 = 16{,}384+64 = 16{,}448\\
\text{fuse }(4\times 64=256\to 64) &: 256\times 64 + 64 = 16{,}384+64 = 16{,}448\\
\text{refine }3\times3 &: 64\times 64\times 9 + 64 = 36{,}864 + 64 = 36{,}928\\
\text{sub-pixel }3\times3 &: 64\times 16\times 9 + 16 = 9{,}216 + 16 = 9{,}232\\[3pt]
\text{total} &: 2{,}112+4{,}160 = 6{,}272\\
 &\phantom{:}\ 6{,}272+10{,}304 = 16{,}576\\
 &\phantom{:}\ 16{,}576+16{,}448 = 33{,}024\\
 &\phantom{:}\ 33{,}024+16{,}448 = 49{,}472\\
 &\phantom{:}\ 49{,}472+36{,}928 = 86{,}400\\
 &\phantom{:}\ 86{,}400+9{,}232 = 95{,}632
\end{align*}
\end{worked}

\begin{sourcenote}[: the head's size, verified]
Building \code{SegFormerDecoder([32, 64, 160, 256], decoder_dim=64, out_channels=1)} on CPU and summing
\code{p.numel()} gives 95,632 \tagC; with \code{out_channels=2} it gives 104,864 \tagC. Script and output:
\file{research/textbook/tools/compute_39-41.py} and \file{compute_39-41_output.txt}. The proposal's ``\textasciitilde
0.1 M parameters'' \tagM{uncertainty-head.md} is this number rounded. For comparison, the
FT04 model it rides on has 5,886,451 parameters with its segmentation head (\cref{ch:shipping}): the head adds
under 2\%.
\end{sourcenote}

\section{The one rule that keeps it honest: hidden pixels only}\label{sec:uncertainty:rule}

The head must judge how hard a place is \emph{from its context, never from the pixel itself}. If the head
could see the pixel, it would learn ``this pixel looks odd, so expect a big error'', and dividing by that $b$
would divide every anomaly away. So the loss is averaged only over pixels that were \emph{hidden from the
encoder} by the masked-autoencoder mask (\cref{ch:masking}), that are valid, and that lie inside the eroded
valid region. The head reads the same encoder features as the reconstruction decoder, and those features were
computed with the hidden tokens removed.

Inference matches training. The product reconstructs a scene in two complementary checkerboard passes
(\cref{ch:inference}), and \code{predict_full_scene_with_scale} runs the same passes with the head switched on,
taking each pixel's $\log b$ from the pass in which it was hidden.

\paragraph{Frozen means frozen.} FT06 trained in a new trainer mode, \code{mode: "scale"}. In the
trainer's constructor (\code{__init__} of \code{FineTuneTrainer}) that mode sets \code{requires_grad=True} only for parameter names starting
with \code{scale_head.}, and builds the optimiser over those alone; the training pass puts the whole model in
evaluation mode and switches only \code{scale_head} back to training mode, so FT04's BatchNorm statistics are not
updated either \tagM{lld/finetune.md}. The reconstruction and the masks are therefore exactly FT04's.

\section{The score: dividing by \texorpdfstring{$b$}{b}, with a floor}\label{sec:uncertainty:score}

The scorer \code{scaled_rm} in \file{app/evaluation/scorers.py} does three things:
\begin{lstlisting}
b = np.exp(log_b)                                                       # (H, W)
b = np.maximum(b, np.percentile(b[valid], scale_floor_pct)) ** scale_power   # floor, then power
score = blurred_mahalanobis((cube - recon) / b[None], valid, spec)       # RM, whitening refitted
\end{lstlisting}
The default floor is the 5th percentile \tagM{scorers.py}. The floor exists because dividing by
a very small $b$ inflates noise wherever the model is usually right (deep water, flat desert): the proposal
listed this as a risk and floored $b$ per scene. \code{scale_power=1} is the design, residual$/b$;
\code{scale_power=0.5} is residual$/\sqrt b$, a softer version tried later.

\begin{worked}[label={we:uncertainty:floor}]{The floor at work}
Suppose a scene's 5th percentile of $b$ is $0.04$. Three pixels have residual $0.20$ and predicted $b$ of
$0.01$, $0.04$ and $0.10$. Compute each scaled residual with the floor, then without.
\begin{align*}
\text{with floor:}\quad & \max(0.01,0.04)=0.04,\quad 0.20/0.04 = 5.0\\
& \max(0.04,0.04)=0.04,\quad 0.20/0.04 = 5.0\\
& \max(0.10,0.04)=0.10,\quad 0.20/0.10 = 2.0\\
\text{without:}\quad & 0.20/0.01 = 20.0,\quad 0.20/0.04 = 5.0,\quad 0.20/0.10=2.0
\end{align*}
Without the floor the first pixel's score is four times larger, purely because the head was very confident
there. The floor caps the reward for confidence.
\end{worked}

\section{Spearman's rank correlation}\label{sec:uncertainty:spearman}

Before any detection test, the proposal set a sanity check: on held-out tiles, $b(p)$ must \emph{rank}
pixels by their actual error, with Spearman correlation at least $0.5$. If the head cannot predict its own
error, the score cannot use it.

\begin{definition}[Rank and Spearman correlation]
The \emph{rank}\index{rank (ordering)} of a value in a list is its position after sorting from smallest to
largest (the repo counts from 0). \emph{Spearman's rank correlation}\index{Spearman rank correlation}
\citep{spearman1904proof} between two equal-length lists $a$ and $e$ is the ordinary (Pearson) correlation of
their ranks:
\begin{equation}\label{eq:uncertainty:spearman}
\rho \;=\; \frac{\sum_i (u_i-\bar u)(v_i-\bar v)}{\sqrt{\sum_i (u_i-\bar u)^2}\,\sqrt{\sum_i(v_i-\bar v)^2}},
\qquad u_i=\operatorname{rank}(a_i),\; v_i=\operatorname{rank}(e_i).
\end{equation}
It is $1$ when the two lists are in the same order, $-1$ when in reverse order, and near $0$ when the orders are
unrelated.
\end{definition}

The repo's \code{spearman} in \file{app/foundation_models/trainers/finetune_trainer.py} does exactly this:
two \code{argsort} calls turn values into ranks, the ranks are centred on their mean, and the result is the dot
product divided by the product of the two vectors' lengths (with ties broken by order, and at most the first
50,000 values). Spearman asks ``does a larger $b$ go with a larger error?'', not ``is $b$ the right size?''.
Doubling every $b$ leaves $\rho$ unchanged.

When there are no ties, \cref{eq:uncertainty:spearman} simplifies to the textbook shortcut
\begin{equation}\label{eq:uncertainty:spearmand}
\rho = 1 - \frac{6\sum_i d_i^2}{n(n^2-1)}, \qquad d_i = u_i - v_i,
\end{equation}
which you can use to check your arithmetic.

\begin{worked}[label={we:uncertainty:spearman5}]{Spearman on five pairs}
Five hidden pixels. The head predicts $b$ and the actual mean error is $e$:
\[
b = (0.12,\ 0.40,\ 0.08,\ 0.25,\ 0.30),\qquad e=(0.10,\ 0.35,\ 0.15,\ 0.20,\ 0.50).
\]
Do every step on paper.

\emph{Step 1, ranks} (0 = smallest). For $b$: the sorted order is $0.08<0.12<0.25<0.30<0.40$, so
$u=(1,4,0,2,3)$. For $e$: $0.10<0.15<0.20<0.35<0.50$, so $v=(0,3,1,2,4)$.

\emph{Step 2, centre.} The mean of the ranks $0,1,2,3,4$ is $10/5=2$ for both lists.
\[
u-\bar u = (-1,\ 2,\ -2,\ 0,\ 1),\qquad v-\bar v=(-2,\ 1,\ -1,\ 0,\ 2).
\]
\emph{Step 3, dot product.}
\begin{align*}
(-1)(-2) + (2)(1) + (-2)(-1) + (0)(0) + (1)(2) &= 2 + 2 + 2 + 0 + 2 = 8.
\end{align*}
\emph{Step 4, lengths.} Both centred vectors are permutations of $(-2,-1,0,1,2)$:
\[
\sqrt{4+1+0+1+4} = \sqrt{10} = 3.162 \text{ each};\qquad 3.162\times 3.162 = 10.
\]
\emph{Step 5.} $\rho = 8/10 = 0.80$.

\emph{Check with \cref{eq:uncertainty:spearmand}.} $d=u-v=(1,1,-1,0,-1)$, so $\sum d^2=1+1+1+0+1=4$, and
\[
\rho = 1-\frac{6\times 4}{5\times(25-1)} = 1-\frac{24}{120} = 1-0.20 = 0.80. \checkmark
\]
The head ranks the errors well but not perfectly: it put pixel 2 (error $0.35$) at the top, where pixel 5
(error $0.50$) belongs.
\end{worked}

\section{The run: FT06}\label{sec:uncertainty:ft06}

FT06's configuration is FT04's (the same mixed batches of EnMAP, PRISMA and AVIRIS-NG at shares 0.4 / 0.35 /
0.25), with \code{mode: "scale"}, \code{scale_decoder_dim: 64}, initialisation from FT04's \code{best.pt}, 3
epochs and learning rate $5\times10^{-4}$ on a cosine schedule
\tagM{FT06-uncertainty/README.md}. The pass criteria, fixed before
launch, were:
\begin{enumerate}
\item the head predicts its own error: validation Spearman$(b,e)\ge 0.5$ on every sensor;
\item detection on the real PRISMA bench against FT04 + RM, both with cloud excluded on FT04's masks: pAUC up
  to 1\% false alarms (\cref{ch:measures}) at least $+0.010$ (0.170 to 0.180), TPR at 0.1\% not lower,
  \emph{and} objects caught at the top 1\% and 0.1\% not lower;
\item bench-v2.1 caught top 1\% not lower by more than 0.5 points.
\end{enumerate}

\subsection{Result 1: the head learned its error}

\begin{table}[htbp]
\centering
\begin{tabular}{lcccc}
\toprule
epoch & EnMAP & PRISMA & AVIRIS-NG & val NLL\\
\midrule
1 & 0.584 & 0.620 & 0.607 & $-1.331$\\
2 & 0.600 & 0.631 & 0.625 & $-1.338$\\
3 & 0.604 & 0.636 & 0.640 & $-1.341$ (best)\\
\bottomrule
\end{tabular}
\caption{FT06 validation Spearman$(b,e)$ per sensor and the head's NLL, per epoch. Source:
\tagM{FT06-uncertainty/README.md} (in \texttt{research/experiments/indradhanu-finetune/}).}
\label{tab:uncertainty:spearman}
\end{table}

\Cref{tab:uncertainty:spearman} shows criterion 1 met on every sensor from the first epoch, and rising every
epoch. The NLL falls each epoch too, and since mode \code{"scale"} ranks epochs by that NLL (``scale by its
NLL'', \file{docs/lld/finetune.md}), \code{best.pt} is epoch 3. Hold on to that: ``best'' here means best at
predicting its own error.

\subsection{Result 2: detection got worse}

\begin{table}[htbp]
\centering\small
\begin{tabular}{lccccc}
\toprule
 & pAUC $\le$1\% & TPR @ 0.1\% & top 1\% & top 0.1\% & ROC-AUC\\
\midrule
FT04 + RM (the bar) & \textbf{0.170} & 0.065 & \textbf{46.0\%} & 15.8\% & \textbf{0.852}\\
FT06 + RM & 0.170 & 0.065 & 46.0\% & 15.8\% & 0.852\\
FT06 + scaled RM, residual$/b$ & 0.125 & 0.064 & 37.8\% & 13.2\% & 0.708\\
FT06 + scaled RM, residual$/\sqrt b$ & 0.163 & \textbf{0.072} & 44.8\% & \textbf{16.6\%} & 0.805\\
\emph{epoch 1} + residual$/\sqrt b$ & 0.176 & 0.079 & 46.1\% & 17.3\% & 0.822\\
\bottomrule
\end{tabular}
\caption{FT06 detection on the real PRISMA bench, cloud excluded on FT04's masks. ``Top'' columns are objects
caught. Source: \tagM{FT06-uncertainty/README.md}.}
\label{tab:uncertainty:detect}
\end{table}

\begin{drift}
FT06 results, run 2026-10-20 (\file{research/experiments/indradhanu-finetune/FT06-uncertainty/}):
\cref{tab:uncertainty:spearman,tab:uncertainty:detect}. bench-v2.1 caught top 1\%: FT04 36.2\%, scaled
33.9\%, sqrt 36.2\%. Subtle tier caught top 1\% with sqrt, cloud excluded: 40.3\% to 37.7\%
\tagM{FT06-uncertainty/eval_report.md}. Set2 pAUC fell 0.128 to 0.097 under residual$/b$.
\end{drift}

Read \cref{tab:uncertainty:detect} row by row.
\begin{itemize}
\item \textbf{FT06 + RM equals FT04 + RM to every digit.} That is the frozen-backbone check working, not a bug:
  the reconstruction is FT04's, so plain RM on it is FT04's score.
\item \textbf{As designed, residual$/b$ loses about a quarter of its pAUC}: $0.170-0.125=0.045$, and
  $0.045/0.170=0.26$. It loses $46.0-37.8=8.2$ points of objects caught at the top 1\%.
\item \textbf{The square-root variant lifts the very top and loses overall}: TPR at 0.1\% goes 0.065 to 0.072,
  but pAUC drops to 0.163.
\item The proposal expected most of any gain on medium and subtle objects. The opposite happened.
\end{itemize}
No row meets criterion 2. FT06 was not adopted.

\section{Why: the head expects large errors exactly where targets are}\label{sec:uncertainty:why}

The diagnosis script \file{b_at_targets.py} asks where the head puts its expected error, on PRISMA sets 1, 2
and 4, for epochs 1 and 3. For each scene it takes the target pixels $T$ and a \emph{ring}\index{ring (around a
target)} around them,
\begin{lstlisting}
ring = binary_dilation(t, iterations=8) & ~binary_dilation(t, iterations=3) & valid
\end{lstlisting}
and reports two numbers: the percentile of $b$ at target pixels against every valid pixel in the scene, and
the \emph{target / ring ratio}
\begin{equation}\label{eq:uncertainty:ratio}
\text{ratio} = \exp\!\big(\operatorname{median}_{T}\log b - \operatorname{median}_{\text{ring}}\log b\big)
 = \frac{\operatorname{median}_T b}{\operatorname{median}_\text{ring} b},
\end{equation}
where the second form holds because $\exp$ is increasing (the median of $\log b$ is the log of the median of
$b$) and $e^{x-y}=e^x/e^y$.

\begin{sourcenote}[: what ``3--8 px'' means]
The FT06 README's observations call the ring ``3-8 px'', and the course repeats it. The code
(\file{FT06-uncertainty/b_at_targets.py}) removes everything within 3 dilation steps and keeps up to 8, so the
ring is pixels at distance 4 to 8 from the target. \code{binary_dilation}'s default structuring element is a
cross, so the distance is the city-block distance $\abs{\Delta x}+\abs{\Delta y}$. The README's evidence line
says this correctly: ``the ring is 4-8 dilation steps, city-block distance''.
\end{sourcenote}

\begin{table}[htbp]
\centering
\begin{tabular}{lccc}
\toprule
 & set1 & set2 & set4\\
\midrule
target / ring $b$ ratio, epoch 1 & 1.05 & 1.12 & 1.04\\
target / ring $b$ ratio, epoch 3 & 1.07 & 1.19 & 1.13\\
\bottomrule
\end{tabular}
\caption{Median $b$ on targets over median $b$ on their own 4--8 px ring. Target pixels also sit at the 87th--90th
percentile of $b$ across each scene. Source: \tagM{FT06-uncertainty/README.md}.}
\label{tab:uncertainty:ratio}
\end{table}

Two separate facts, and both hurt.
\begin{itemize}
\item \textbf{Across the scene}, targets sit on built-up surfaces, which are hard to reconstruct, so the head
  rightly expects large errors there (the 87th--90th percentile).
\item \textbf{Locally}, $b$ on the target is higher than on its own ring. That should not happen if the head only
  judged ``how hard is this kind of place'': the ring is the same kind of place.
\end{itemize}

\subsection{Why the hidden-pixel rule did not stop it: halo arithmetic}

The rule hides the pixel from the encoder. But the encoder hides whole $4\times4$-pixel tokens (stage 1 is a
non-overlapping kernel-4, stride-4 convolution, \cref{ch:patchembed}), and a target is not confined to its
token. Its \emph{PSF halo}\index{PSF halo} (the sensor's blur that spreads some of a target's light into
neighbouring pixels, \cref{ch:benchmarks}) and any part of a target larger than one token can sit in the
\emph{neighbouring} tokens, which are visible about 35\% of the time.

\begin{worked}[label={we:uncertainty:halo}]{How often a halo leaks out of its token}
Take a one-pixel target whose light reaches $h$ pixels on either side. Its footprint along one axis is
$2h+1$ pixels wide. Assume (an illustrative assumption \tagA) the target's position within its 4-pixel token is
equally likely to be any of the 4 offsets $0,1,2,3$. The footprint stays inside the token along that axis only
if it starts at offset $\ge 0$ and ends at offset $\le 3$.
\begin{align*}
h=0:&\ \text{width }1;\ \text{fits at offsets }0,1,2,3 \to 4\text{ of }4 = 1\\
h=1:&\ \text{width }3;\ \text{the target must be at offset 1 or 2} \to 2\text{ of }4 = 1/2\\
h=2:&\ \text{width }5 > 4 \to 0\text{ of }4 = 0
\end{align*}
Both axes must fit, so for $h=1$ the footprint is inside one token with probability
$\tfrac12\times\tfrac12=\tfrac14$: three times in four, some of the target's light lands in a neighbouring
token. For $h=2$ it always does. So for almost every target, a visible neighbour carries a trace of it.
\end{worked}

The head learns ``something odd is next to here, expect a big error''. It is right about that: the odd
spectrum next door genuinely predicts a large reconstruction error. And dividing by $b$ then explains the
anomaly away.

\begin{keypoint}
A head trained to predict its own error will predict a large error on an anomaly whenever it can see any
trace of it, and then the anomaly is no longer surprising.
\end{keypoint}

\begin{worked}[label={we:uncertainty:divide}]{Dividing by $b$, by hand}
Use set2's epoch-3 ratio, $1.19$. Keep to one whitened band, ignore that the whitening is refitted (it rescales
every pixel alike, so it does not change the order), and score by squared distance from a ring mean of $0$
(RM's score is a squared distance). Three pixels: a target with residual $3.0$ and $b=1.19$; a pixel in its ring
with residual $1.0$ and $b=1.00$; a flat-field pixel elsewhere with residual $2.6$ and $b=1.00$. The residuals
are illustrative \tagA; only the ratio $1.19$ is measured.

\emph{Plain RM:}
\begin{align*}
\text{target} &: 3.0^2 = 9.00\\
\text{ring pixel} &: 1.0^2 = 1.00\\
\text{flat field} &: 2.6^2 = 6.76 \qquad\text{(target ranks first)}
\end{align*}
\emph{Residual$/b$:}
\begin{align*}
\text{target} &: 3.0/1.19 = 2.521,\quad 2.521^2 = 6.36\\
\text{ring pixel} &: 1.0/1.00 = 1.000,\quad 1.000^2 = 1.00\\
\text{flat field} &: 2.6/1.00 = 2.600,\quad 2.600^2 = 6.76 \qquad\text{(the flat field now outranks the target)}
\end{align*}
The target's score fell from $9.00$ to $6.36$: divided by $1.19\times1.19=1.416$, about 29\% lower. It fell
below an ordinary bright pixel elsewhere, which is how objects drop out of the top 1\%.

\emph{Residual$/\sqrt b$:} $\sqrt{1.19}=1.091$, so $3.0/1.091=2.750$, and $2.750^2=7.56$, still above 6.76,
which is why the square-root variant hurts less.
\end{worked}

\begin{figure}[htbp]
\centering
\begin{tikzpicture}
\begin{axis}[ybar,width=0.85\linewidth,height=5.5cm,bar width=12pt,ymin=0,ymax=10,
  symbolic x coords={plain RM,residual/sqrt b,residual/b},xtick=data,ylabel={squared score},
  legend style={at={(0.5,1.02)},anchor=south,legend columns=3},enlarge x limits=0.25,
  nodes near coords,every node near coord/.append style={font=\scriptsize}]
\addplot[fill=loss!70] coordinates {(plain RM,9.00) (residual/sqrt b,7.56) (residual/b,6.36)};
\addplot[fill=maskgrey!60] coordinates {(plain RM,6.76) (residual/sqrt b,6.76) (residual/b,6.76)};
\addplot[fill=spec!50] coordinates {(plain RM,1.00) (residual/sqrt b,1.00) (residual/b,1.00)};
\legend{target,flat field,ring pixel}
\end{axis}
\end{tikzpicture}
\caption{The three pixels of \cref{we:uncertainty:divide}, drawn to scale. Dividing by $b$ with a target/ring
ratio of 1.19 drops the target below an ordinary flat-field pixel; the square root keeps it a little above.}
\label{fig:uncertainty:divide}
\end{figure}

\subsection{Why epoch 1 looked better than epoch 3}

The ratios in \cref{tab:uncertainty:ratio} grew with training: set1 1.05 to 1.07, set2 1.12 to 1.19, set4 1.04
to 1.13. As the head got better at its own job (Spearman up every epoch), it also learned more of the
``something odd next door'' signal, because that signal genuinely predicts large errors. More of it means more
of each target divided away. So epoch 1 with the square root scored pAUC 0.176, and epoch 3, the
\code{best.pt} by NLL, fell below FT04.

\begin{keypoint}
Predicting your own error well is not the same goal as detection: the checkpoint chosen by the head's own
loss was the worst one for the product.
\end{keypoint}

\subsection{The selection-risk caveat}

Do not read the epoch-1 row as a gain to ship. The square-root variant was \emph{chosen after seeing epoch-1
results on the real bench}, and then judged on the same bench. When you try variants and keep the one that
looks best on the data that will also grade it, part of the improvement is luck in that data, and you cannot
tell how much. This is \emph{selection risk}\index{selection risk}. Even at face value, epoch 1's pAUC of 0.176
is below the 0.180 bar. \Cref{ch:bootstrap} later put a paired bootstrap interval on the epoch-3 square-root
row: $-0.0067\ [-0.0174, +0.0037]$ \tagM{bench-resolution/README.md}, a tie.

\section{What might rescue it, and what it taught}\label{sec:uncertainty:rescue}

The FT06 README lists two ideas, both untried and both costing a retrain:
\begin{enumerate}
\item hide a wider context when training the head, for example $3\times3$ tokens around each predicted one, so
  the head never sees a target's halo or its other parts;
\item predict $b$ from the reconstruction branch's features of a heavily masked pass only.
\end{enumerate}
The README recommends neither before a cheaper idea. The halo problem did not go away: it is the motivation
for the block masking of \cref{ch:decoderwork}.

A practical note from the run: the evaluation took about 2.3 hours rather than the expected 1, because two
models plus the scale pass filled the laptop GPU's 20\,GB \tagM{FT06-uncertainty/README.md}.

\begin{confusions}
\confusion{Spearman met its bar, so the head works.}{It works at its own goal. Criterion 1 was a sanity check:
a head that cannot predict its error is useless. Passing it was necessary, not sufficient; detection is the
criterion that counts.}
\confusion{The hidden-pixel rule should have prevented this.}{It prevents the head from seeing the pixel. It
does not stop the head from seeing neighbouring tokens, which carry the target's halo and any part larger than
one $4\times4$ token (\cref{we:uncertainty:halo}).}
\confusion{FT06 + RM matching FT04 to every digit means the evaluation is broken.}{It is the check that the
backbone stayed frozen. Different numbers there would have been the bug.}
\confusion{The 1.19 is a percentile.}{No. It is a ratio of medians, target $b$ over ring $b$. The 87th--90th
percentile is a separate measure, of target $b$ against the whole scene.}
\confusion{The Laplace NLL has a log 2 that the code forgot.}{It was dropped on purpose: a constant does not
change where the minimum is, so it does not change any gradient.}
\confusion{Why predict log b instead of b?}{A convolution can output negative numbers; $e^{s}$ is always
positive. It also turns the $1/b^2$ gradient near zero into a gentle one (\cref{we:uncertainty:chain}).}
\confusion{This is the same as RM's whitening.}{RM divides by a per-scene spread along each error
\emph{direction}. $b(p)$ is a per-\emph{pixel} scale. The design keeps per-band shape with RM on purpose.}
\confusion{best.pt is the best model.}{It is the epoch with the lowest scale NLL, which mode
\texttt{"scale"} uses to rank epochs. Here that was the worst epoch for detection.}
\end{confusions}

\begin{checkyourself}
\item\label{ex:uncertainty:1} For $\abs{r}=0.5$, compute $\ell=\abs{r}/b+\ln b$ at $b=0.25$, $0.5$ and $1.0$.
Which is lowest? Then compute $d\ell/db$ at each of the three.
\item\label{ex:uncertainty:2} Five pixels have $b=(0.3, 0.1, 0.5, 0.2, 0.4)$ and errors
$e=(0.25,\allowbreak 0.05,\allowbreak 0.30,\allowbreak 0.40,\allowbreak 0.45)$. Compute Spearman's $\rho$ in two ways: with \cref{eq:uncertainty:spearman},
and with \cref{eq:uncertainty:spearmand}.
\item\label{ex:uncertainty:3} What would the head's parameter count be with decoder width 32 (and still one output
channel)? Count every layer.
\item\label{ex:uncertainty:4} Redo \cref{we:uncertainty:divide} with set1's epoch-3 ratio, 1.07. What is the
target's squared score under residual$/b$, and does it still outrank the flat field at 6.76? By what factor
does residual$/\sqrt b$ shrink the squared score with ratio 1.19?
\item\label{ex:uncertainty:5} A Gaussian error model gives the NLL $r^2/(2\sigma^2)+\log\sigma$ (constants
dropped). Find the best $\sigma$ for one error $r$, and for $N$ errors sharing one $\sigma$. Compare with the
Laplace result.
\item\label{ex:uncertainty:6} Spearman rose every epoch while detection fell. Name the mechanism, and say why
the epoch-1 square-root row is not evidence of a gain.
\end{checkyourself}

\nextchapter{ch:decoderwork}
